\section{讨论：有效性边界、识别任务与理论完成度}

\subsection{有效性边界}

本文框架的严格内容在于：切投影与扫描读出可将静态高维离散结构映射为三维准周期几何与时间序列；有限分辨率读出诱导等价类与统计；稳定扇区提供“对象化”结构的协议定义。而将稳定类型进一步识别为特定粒子谱、将缺陷通道识别为特定相互作用场，属于识别层工作，必须通过明确的数据拟合与可证伪中间命题推进。

\subsection{识别任务}

理论完成需要解决的核心识别任务包括：
\begin{enumerate}
  \item 在具体实验数据上构造可逆的“读出—类型”重编码管线，验证退化指纹与准周期谱；
  \item 给出跨尺度的分辨率提升规则（$m\to m+1$）并验证其与观测统计的分层一致性；
  \item 给出缺陷通道的可计算连接结构，使其可对应到可观测的相互作用补偿项或约束一致性条件。
\end{enumerate}

